\section{Tail Leakage}
\label{sec:tail-leakage}

Let \(p_t(v)\) be the next-token distribution at position \(t\) during
greedy generation, and let \(V_k(p_t)\) contain its \(k\) most probable
tokens. We define the retained mass
\begin{equation}
    m_t^{(k)}=\sum_{v\in V_k(p_t)}p_t(v),
    \qquad
    \ell_t^{(k)}=1-m_t^{(k)},
    \label{eq:retained-mass}
\end{equation}
where \(\ell_t^{(k)}\) is the mass that leaks outside the leading set. We use
\(k=5\) throughout the main experiments and write \(m_t\) for
\(m_t^{(5)}\).

\subsection{Why isolate leaked mass?}

Token entropy measures the total spread of \(p_t\), but this spread has
several sources. Let \(\hat p_t^{\,\mathrm{in}}\) and
\(\hat p_t^{\,\mathrm{out}}\) denote the distributions obtained by
renormalizing the tokens inside and outside \(V_5(p_t)\), respectively. The
entropy chain rule gives
\begin{align}
H(p_t)={}&-m_t\log m_t-(1-m_t)\log(1-m_t) \nonumber\\
         &+m_t H(\hat p_t^{\,\mathrm{in}})
          +(1-m_t)H(\hat p_t^{\,\mathrm{out}}).
\label{eq:entropy-decomposition}
\end{align}
The first term records how much mass crosses the top-five boundary. The other
terms record competition among the leading candidates and the arrangement of
probability within the tail. For short-answer QA, competition among leading
tokens can reflect casing, aliases, or tokenization rather than factual
uncertainty. Conversely, the internal shape of a very large tail may change
without materially changing the model's leading hypothesis. Equation
\ref{eq:entropy-decomposition} motivates testing the boundary-crossing mass
separately instead of compressing all three effects into entropy.

\subsection{Token- and answer-level score}

We convert retained mass into a monotonic risk score:
\begin{equation}
    \tl_t=-\log m_t,
    \qquad
    \tl(a)=\sum_{t\in\mathcal{T}(a)}\tl_t
          =-\log\prod_{t\in\mathcal{T}(a)}m_t ,
\label{eq:tail-leakage}
\end{equation}
where \(\mathcal{T}(a)\) contains the generated answer tokens and excludes
EOS. The negative logarithm has two useful properties. First, it increases
strictly as mass leaves the top-five set; unlike binary entropy, it does not
decrease once leaked mass exceeds \(0.5\). Second, it is additive across
positions, allowing small amounts of leakage distributed over a multi-token
answer to accumulate.

\tailleak{} requires only the top-\(k\) probabilities already available during
a greedy pass. It is therefore cheaper than sample-disagreement or semantic
entropy estimates that require multiple generations. It is a ranking score,
not a calibrated probability that an answer is wrong. In particular, if a
model concentrates almost all probability on a false continuation,
Equation~\ref{eq:tail-leakage} provides no warning.

\section{Adaptive Retrieval with \method{}}
\label{sec:seek}

We use \method{} to test whether \tailleak{} is useful beyond offline error
detection. Given a question \(q\), a language model \(\mathcal{M}\), and a
passage corpus \(\mathcal{P}\), \method{} preserves two answer candidates until
late routing:

\begin{figure}[t]
\centering
\small
\setlength{\fboxsep}{6pt}
\fbox{\parbox{0.93\linewidth}{
\textbf{Direct path:}
\(q \rightarrow\) retrieve with \(q\) \(\rightarrow a^{\mathrm{dir}}\)
\(\quad\)\(\Vert\)\(\quad\)
\textbf{Decomposition path:}
\(q\rightarrow(s_1,\ldots,s_n)\rightarrow\) gate retrieval at each \(s_i\)
\(\rightarrow a^{\mathrm{dec}}\).
\par\medskip
\centering
\(a^{\mathrm{dir}},a^{\mathrm{dec}}\rightarrow\)
transparent late routing \(\rightarrow\hat a\)
}}
\caption{The two candidate paths in \method{}. Decomposition is retained as a
retrieval strategy rather than an irreversible early decision.}
\label{fig:seek-overview}
\end{figure}

\paragraph{Executable decomposition.}
A LoRA-adapted decomposer maps \(q\) to ordered, knowledge-graph-style
sub-questions. A step has the form
\(\texttt{entity >> relation}\); a placeholder such as \(\texttt{\#1}\)
refers to the answer of an earlier step. For example,
\(\texttt{Inception >> director}\) followed by
\(\texttt{\#1 >> place of birth}\) becomes a directly executable second query
after the first answer is substituted. One-step reformulations are allowed,
which prevents the representation from forcing a multi-hop chain onto
single-hop questions.

\paragraph{Step-level retrieval gate.}
For each resolved sub-question \(s_i\), the model first generates a tentative
answer \(\tilde a_i\) without retrieved evidence and computes
\(\tl(\tilde a_i)\). Retrieval is invoked when
\begin{equation}
    \pi(s_i)=\mathbb{I}\!\left[\tl(\tilde a_i)\geq\tau\right].
\label{eq:gate}
\end{equation}
If the gate fires, BM25 returns passages and the model re-answers \(s_i\);
otherwise \(\tilde a_i\) is retained. Thus the signal operates at a natural
reasoning boundary rather than at an arbitrary sentence or token position.

\paragraph{Context preservation and synthesis.}
Short entity-centered sub-questions are useful retrieval queries but can omit
comparison targets or disambiguating constraints. When a step retrieves,
\method{} therefore deduplicates the union of passages retrieved from \(s_i\)
and from the original question \(q\). After all placeholders are resolved, a
base-model synthesis pass reads \(q\) and the ordered \((s_i,a_i)\) pairs to
produce \(a^{\mathrm{dec}}\). The direct candidate \(a^{\mathrm{dir}}\)
answers \(q\) from original-question evidence.

\paragraph{Late routing.}
A deterministic verifier defaults to the decomposition candidate and switches
to \(a^{\mathrm{dir}}\) only when direct-answer \tailleak{} is low and the
decomposition path shows instability, such as high post-synthesis leakage,
malformed placeholders, or an incompatible answer type. The verifier does not
regenerate an answer or train a policy. We include it to preserve complementary
candidates, but do not treat it as the principal source of accuracy:
Section~\ref{sec:routing-results} shows that its net effect is small. Exact
rules and diagnostic counts appear in Appendix~\ref{app:details}.
